\section*{Bab \ref{ch:b3:asymmetric-synthesis} --- Sintesis Asimetris}

\begin{solution}{exo:b3:asymmetric-synthesis:1}
ee 80\,\%: $\mathrm{er} = 1.8/0.2 = 9$ (90\,:\,10). er 99\,:\,1: $\mathrm{ee} = 98/100 = 98$\,\%. ee 50\,\%: 75\,\% dan 25\,\%.
\end{solution}

\begin{solution}{exo:b3:asymmetric-synthesis:2}
$(1840 - 60)/(1840 + 60) = 0.937$: ee 93.7\,\%, yaitu er $1840/60 = 31$.
\end{solution}

\begin{solution}{exo:b3:asymmetric-synthesis:3}
Prioritas O $>$ etil $>$ H. Jika digambar dengan O di atas, etil di kiri bawah, dan H di kanan bawah,
muka depannya adalah \textit{Si} dan muka belakangnya \textit{Re}. Gugus metil yang ditambahkan dari \omterm{def:b3:asymmetric-synthesis:faces}{muka Re}, dengan
prioritas OH $>$ etil $>$ metil $>$ H dalam produk, memberikan (\textit{R})-butan-2-ol.
\end{solution}

\begin{solution}{exo:b3:asymmetric-synthesis:4}
Etanol: \omterm{def:b3:asymmetric-synthesis:topicity}{enantiotopik} (dipertukarkan oleh \omterm{def:b3:point-groups:mirror}{bidang cermin} molekul; geseran
sama dalam pelarut akiral). Butan-2-ol: \omterm{def:b3:asymmetric-synthesis:topicity}{diastereotopik} (bersebelahan dengan pusat stereo;
geseran berbeda).
\end{solution}

\begin{solution}{exo:b3:asymmetric-synthesis:5}
er 199: $RT\ln 199 = \qty{13.1}{kJ/mol}$ pada \qty{298}{K}, \qty{8.6}{kJ/mol} pada \qty{195}{K}. ee 80\,\% pada \qty{25}{\celsius} adalah er 9,
$\Delta\Delta G^\ddagger = RT\ln 9 = \qty{5.45}{kJ/mol}$; pada \qty{-78}{\celsius}, $\mathrm{er} = \eu^{5450/(8.314 \times 195.15)} = 28.7$, ee 93\,\%.
\end{solution}

\begin{solution}{exo:b3:asymmetric-synthesis:6}
\omterm{def:b3:asymmetric-synthesis:ee}{Kemurnian optis} $12.0/15.0 = 80$\,\%: 90\,\% (\textit{S}) dan 10\,\% (\textit{R}). Rotasi yang kecil, pengotor,
efek konsentrasi dan pelarut, serta hubungan rotasi--komposisi yang tidak linear
membuatnya tidak andal; kromatografi memisahkan dan menghitung
enantiomer-enantiomer itu secara langsung.
\end{solution}

\begin{solution}{exo:b3:asymmetric-synthesis:7}
$s = \ln(0.40 \times 0.10)/\ln(0.40 \times 1.90) = -3.22/-0.274 = 12$.
\end{solution}

\begin{solution}{exo:b3:asymmetric-synthesis:8}
Dengan L-($+$)-dietil tartrat, (2\textit{S},3\textit{S})-2,3-epoksiheksan-1-ol; dengan D-($-$)-dietil tartrat,
enantiomernya, (2\textit{R},3\textit{R}).
\end{solution}

\begin{solution}{exo:b3:asymmetric-synthesis:9}
L = fenil, M = metil, S = H. Untuk aldehida (\textit{S}), konformer reaktifnya memiliki fenil
yang tegak lurus terhadap bidang karbonil dan hidrogen di samping lintasan
nukleofil, yang menyerang secara anti terhadap fenil. Penentuan konfigurasinya memberikan
(2\textit{S},3\textit{S})-3-fenilbutan-2-ol sebagai diastereomer utama.
\end{solution}

\begin{solution}{exo:b3:asymmetric-synthesis:10}
Pada awalnya kedua enantiomer hadir dalam jumlah yang sama, sehingga produk-produknya
terbentuk dengan perbandingan $k_{\mathrm{fast}} : k_{\mathrm{slow}} = s$: $\mathrm{ee}_p = (s - 1)/(s + 1)$, 0.905 untuk $s = 20$. Selama reaksi berlangsung,
enantiomer yang cepat terkuras dan ee produk turun, sedangkan ee substrat naik
menuju 1: substrat selalu dapat dimurnikan dengan melanjutkan reaksi; produk
tidak dapat.
\end{solution}

\begin{solution}{exo:b3:asymmetric-synthesis:11}
Dengan rasemisasi yang cepat: rendemen sampai 100\,\%, $\mathrm{er} = 99$, ee 98\,\%. Tanpa rasemisasi, pada konversi
50\,\%: rendemen paling banyak 50\,\%, dan ee produk 93\,\% (ee substratnya
juga 93\,\%).
\end{solution}

\begin{solution}{exo:b3:asymmetric-synthesis:12}
Jika kedua kompleks saling berubah lebih cepat daripada laju reaksinya masing-masing, rasio produknya adalah
$k_B[\mathrm B]/(k_A[\mathrm A]) = 600/10 = 60$: kompleks minor B memberikan 98\,\% produk (ee 97\,\%).
Selektivitas ditentukan oleh selisih antara kedua keadaan transisi, bukan oleh
\omterm{def:b3:partition-functions:boltzmann}{populasi} zat-zat antara.
\end{solution}

\begin{solution}{pb:b3:asymmetric-synthesis:1}
\textbf{1.} Karbon alkena yang membawa nitrogen memiliki tiga substituen yang berbeda
dan dua muka yang berbeda; hidrogen yang ditambahkan padanya menjadikan karbon-$\alpha$ asam
amino itu sebuah pusat stereo.

\textbf{2.} Rasemat: kedua muka diserang dengan laju yang sama.

\textbf{3.} Oksigen karbonilnya mengikat rodium bersama ikatan C=C: substrat dipegang
sebagai kelat, dalam geometri tetap yang dapat dibedakan oleh ligan.

\textbf{4.} Difosfina itu membangun kantong kiral di sekeliling logam (bersimetri $C_2$): kedua cara
mengikat substrat, melalui muka yang satu atau yang lain, bersifat diastereomer, dengan energi
dan kereaktifan yang berbeda.

\textbf{5.} Tidak: kompleks minor bereaksi jauh lebih cepat dengan hidrogen (Curtin--Hammett).

\textbf{6.} Logam berada pada satu muka alkena yang terikat, dan hidrogen-hidrogennya
dipindahkan dari logam.

\textbf{7.} $\mathrm{er} = 1.95/0.05 = 39$.

\textbf{8.} $\Delta\Delta G^\ddagger = RT\ln 39 = 8.314 \times 298.15 \times 3.664 = \qty{9.1}{kJ/mol}$.

\textbf{9.} $\mathrm{er} = \eu^{9082/(8.314 \times 195.15)} = 270$: ee 99.3\,\%.

\textbf{10.} $\mathrm{er} = \eu^{11082/(8.314 \times 298.15)} = 87$: ee 97.7\,\%.

\textbf{11.} Kurang dari sebuah ikatan hidrogen yang khas: perbedaan kecil dalam kontak
sterik atau satu interaksi lemah antara substrat dan ligan menentukan
hasilnya, dan itulah sebabnya katalis semacam itu sulit dirancang dari prinsip-prinsip dasar.

\textbf{12.} Lajunya turun, kelarutan hidrogen dan perilaku katalis berubah
pada suhu rendah, dan mendinginkan reaktor yang besar memerlukan energi.

\textbf{13.} 97.5 g enantiomer (\textit{S}) dan 2.5 g enantiomer (\textit{R}).

\textbf{14.} $97.5 - 5.0 = \qty{92.5}{g}$ kristal, yang praktis murni enantiomer.

\textbf{15.} $92.5/100 = 92.5$\,\% dari produk mentah (95\,\% dari enantiomer (\textit{S})).

\textbf{16.} Kristal-kristal mengambil enantiomer yang berlebih, sedangkan bagian rasematnya tetap dalam
larutan bersama enantiomer minor; dari rasemat tidak ada kelebihan yang dapat dikumpulkan, dan
kristalisasi memberikan kristal rasemat.

\textbf{17.} Dengan HPLC kiral terhadap acuan rasemat.

\textbf{18.} 50\,\%: hanya enantiomer yang lambat yang bertahan.

\textbf{19.} $s = \ln(0.45 \times 0.05)/\ln(0.45 \times 1.95) = -3.79/-0.131 = 29$.

\textbf{20.} 45\,\% rasemat (88\,\% enantiomer yang dikehendaki yang ada pada awalnya,
dengan ee 95\,\%).

\textbf{21.} Dengan rasemisasi substrat yang cepat, rendemennya dapat mendekati 100\,\%.

\textbf{22.} Hidrogenasi itu mengubah seluruh prekursor akiral yang murah menjadi enantiomer yang tepat,
dengan katalis yang dipakai dalam jumlah sangat kecil dan tanpa enantiomer yang terbuang.

\textbf{23.} Rodium beracun dan residunya dalam obat dibatasi sampai beberapa bagian per
juta; residu itu diukur dengan spektrometri massa plasma gandeng induktif.

\textbf{24.} Ee 95\,\% pada \qty{25}{\celsius} bersesuaian dengan $\Delta\Delta G^\ddagger = RT\ln 39 \approx \qty{9.1}{kJ/mol}$.
\end{solution}
